\chapter{Conclusiones y trabajo a futuro}
\label{chapter:conclu}

En este trabajo llevamos a cabo la implementación completa del modelo de \citet{echeveste2020cortical} dentro del framework Scikit-NeuroMSI, estableciendo las bases técnicas necesarias para futuras investigaciones que combinen dinámicas corticales realistas con integración multisensorial. Esta implementación representa el primer paso hacia una visión más amplia: el desarrollo de modelos que unifiquen la arquitectura multimodal del enfoque de Cuppini con los mecanismos de muestreo estocástico del modelo de Echeveste.

La reimplementación del código original, que se encontraba en OCaml, requirió una traducción completa a Python utilizando JAX para diferenciación automática. Este proceso, aunque inicialmente inesperado, resultó en una implementación moderna compatible con el ecosistema actual de herramientas de neurociencia computacional, con soporte transparente para aceleración en \acrshort{GPU} y una arquitectura modular que facilita futuras extensiones.

La validación de la implementación se realizó en múltiples niveles complementarios. La verificación función por función demostró equivalencia numérica exacta con el código original, alcanzando errores del orden de la precisión de punto flotante ($\sim 10^{-15}$). Las pruebas del procedimiento de entrenamiento confirmaron que los parámetros optimizados convergen hacia valores comparables a los de referencia, con múltiples corridas independientes mostrando consistencia en los resultados. Las simulaciones replicaron las figuras centrales del artículo original, verificando que la red implementada reproduce correctamente las propiedades de inferencia probabilística mediante muestreo y las dinámicas temporales características, incluyendo oscilaciones gamma cuya frecuencia aumenta con el contraste del estímulo.

Como contribuciones adicionales al framework, se desarrollaron dos paquetes nuevos: \texttt{data} que proporciona infraestructura para gestionar parámetros pre-entrenados y componentes de modelos generativos, y \texttt{generative} que introduce por primera vez modelos generativos como el \acrshort{GSM} al framework. Este último módulo fue diseñado con extensibilidad en mente, permitiendo incorporar futuros modelos generativos siguiendo la misma interfaz.


% Es importante reconocer las limitaciones del presente trabajo. La implementación se centró exclusivamente en el modelo unimodal de Echeveste, sin abordar la fusión con el modelo de Cuppini ni la implementación de inferencia causal multisensorial. Estas extensiones, que constituían la motivación inicial del proyecto, requieren resolver desafíos técnicos y conceptuales adicionales que exceden el alcance de esta tesis.
% TODO: Acá tendrías que ser más específico para mencionar por qué excede el alcance, si inicialmente fue tu objetivo. O sea, sabemos que te faltó tiempo, recursos, etc, pero debes dejarlo claro. Si no, suena a que tiene problemas que están fuera de tu alcance. Vale la pena ser más detallado para explicar a qué te refieres con desafíos técnios y conceptuales. Si no, suena a chamuyo.

Es importante reconocer las limitaciones del presente trabajo. La implementación se centró exclusivamente en el modelo unimodal de Echeveste, sin abordar la fusión con el modelo de Cuppini ni la implementación de inferencia causal multisensorial. Estas extensiones, que constituían la motivación inicial del proyecto, quedaron fuera del alcance de esta tesis porque la implementación del modelo de Echeveste resultó ser un trabajo considerablemente mayor al anticipado. Lo que inicialmente se planteó como una adaptación relativamente directa se transformó en un desarrollo de gran tamaño, el descubrimiento de que el código de entrenamiento original estaba escrito en OCaml requirió una reimplementación completa desde cero, la integración con el framework demandó resolver incompatibilidades de diseño no previstas, y cada componente del modelo (el \acrshort{GSM}, las dinámicas de la red, el procedimiento de entrenamiento, la visualización de resultados) presentó problemas técnicos propios que fueron expandiendo progresivamente el volumen de trabajo. El resultado fue una implementación más extensa y compleja de lo planificado, pero también más completa y rigurosa. Todo esto teniendo en cuenta que la implementación validada que aquí presentamos constituye una condición necesaria para cualquier desarrollo futuro en esa dirección.


% Como resultado, la implementación del modelo de Echeveste convierte a Scikit-NeuroMSI en una plataforma más completa para el estudio de modelos neurocomputacionales, incorporando por primera vez un modelo con dinámicas corticales basadas en muestreo estocástico. Esta contribución no solo beneficia a quienes deseen utilizar específicamente este modelo, sino que establece patrones y herramientas que pueden ser aprovechados para la implementación de modelos similares en el futuro.
% TODO: ¿Qué modelos por ejemplo? Acá es donde debes citar, para que quede claro a qué literatura científica vas a aportar.
Como resultado, la implementación del modelo de Echeveste convierte a Scikit-NeuroMSI en una plataforma más completa para el estudio de modelos neurocomputacionales, incorporando por primera vez un modelo con dinámicas corticales basadas en muestreo estocástico. Esta contribución no solo beneficia a quienes deseen utilizar específicamente este modelo, sino que establece patrones y herramientas que pueden ser aprovechados para la implementación de modelos similares. Por ejemplo, el trabajo de \citet{ursino2017multisensory} demostró que las sinapsis cross-modales del modelo de Cuppini pueden entrenarse mediante aprendizaje hebbiano para reflejar conocimiento a priori sobre la co-ocurrencia de estímulos; una extensión natural sería incorporar este mecanismo de maduración sináptica dentro del framework. De manera análoga, el modelo de \citet{haefner2016perceptual}, que extiende el paradigma de muestreo neuronal a la toma de decisiones perceptuales, representa otro candidato para implementación en Scikit-NeuroMSI aprovechando la infraestructura que aquí desarrollamos para modelos basados en muestreo estocástico.


Más allá de su contribución técnica específica, este trabajo se enmarca en una perspectiva más amplia sobre el rol de los modelos biológicamente plausibles tanto en la comprensión del cerebro como en el futuro de la neurociencia computacional. Actualmente existe una brecha profunda entre las redes neuronales artificiales, optimizadas exclusivamente para maximizar rendimiento en tareas específicas, y los circuitos neuronales biológicos que exhiben propiedades como estocasticidad, dinámicas oscilatorias, y estricto balance entre excitación e inhibición. El modelo de Echeveste demuestra que estas propiedades biológicas no son meros epifenómenos evolutivos sino que cumplen roles computacionales precisos: las oscilaciones gamma que emerge de la red entrenada aceleran la exploración del espacio de estados \citep{echeveste2020cortical}, y la variabilidad neuronal codifica la incertidumbre de la distribución posterior del modelo generativo \citep{echeveste2020cortical}. Estos resultados se suman a evidencia previa mostrando que la actividad espontánea cortical refleja propiedades estadísticas del modelo interno del entorno \citep{berkes2011spontaneous}, y que la variabilidad neuronal modulada por el estímulo es consistente con un mecanismo de muestreo probabilístico \citep{banyai2019stimulus}.

\section{Trabajo a futuro}

% La dirección más prometedora para trabajo futuro es la integración del modelo de Echeveste con el modelo de Cuppini ya implementado en Scikit-NeuroMSI. Esta fusión permitiría estudiar cómo fenómenos de integración multisensorial, como el efecto ventrílocuo, emergen de circuitos neuronales que implementan inferencia probabilística mediante muestreo estocástico. Tal modelo unificado tendría implicaciones tanto teóricas, demostrando que los principios de muestreo neuronal pueden extenderse al dominio multisensorial, como prácticas, permitiendo conectar predicciones del modelo con observaciones electrofisiológicas reales de integración audiovisual.

% TODO: ¿Realmente sería la primera vez que se hace algo así? ¿El modelo de Ursino no es algo parecido? Debes ser más preciso en qué aspecto reside la innovación. Sugiero que puedas mencionar también la literatura en neurociencias sobre ondas EEG e integración multisensorial. Ahí vas a encontrar detalles sobre las ondas beta, gamma, que pueden ser asociados al paradigma del ventrículo espacial. Mira textos de Romei por ejemplo, o en mis papers las referencias. Tu innovación sería esa. Puedes modelar directamente las ondas cerebrales asociadas a la integración multisensorial! Eso sí no se ha hecho antes que yo sepa.

La dirección más prometedora para trabajo futuro es la integración del modelo de Echeveste con el modelo de Cuppini ya implementado en Scikit-NeuroMSI. Esta fusión permitiría estudiar cómo fenómenos de integración multisensorial, como el efecto ventrílocuo, emergen de circuitos neuronales que implementan inferencia probabilística mediante muestreo estocástico. Trabajos previos del grupo de Ursino ya han avanzado en direcciones relevantes, \citet{ursino2009recognition} exploraron el rol de la sincronización gamma en el reconocimiento de objetos mediante osciladores neuronales, y \citet{ursino2017multisensory} demostraron que las sinapsis cross-modales del modelo de Cuppini pueden aprenderse mediante maduración hebbiana, incorporando conocimiento a priori sobre la co-ocurrencia de estímulos. La fusión que proponemos busca llevar estos avances un paso más allá, incorporando dinámicas corticales basadas en muestreo estocástico con balance excitación-inhibición al dominio multisensorial. A diferencia de los modelos deterministas de tasa existentes, un modelo así generaría señales de \acrshort{LFP} simuladas con oscilaciones emergentes en bandas fisiológicamente relevantes (gamma, beta, theta), lo que abriría la posibilidad de modelar directamente las dinámicas oscilatorias cerebrales asociadas a la integración multisensorial y contrastarlas con registros electrofisiológicos reales de tareas como la ventriloquía espacial.

Aunque nuestra implementación reproduce y verifica varias propiedades reportadas en el artículo original, una línea complementaria de trabajo futuro consiste en extender la validación experimental del modelo. Este tipo de validaciones tiene precedentes claros, donde por ejemplo, \citet{echeveste2020cortical} entrenaron redes de control que diferían en un único aspecto (sin modulación de covarianza, sin principio de Dale, sin penalización de slowness) para demostrar la necesidad de cada componente, y \citet{banyai2019stimulus} realizaron análisis detallados de cómo la complejidad del estímulo modula las correlaciones de respuesta en \acrshort{V1}. Replicar estas comparaciones con nuestra implementación fortalecería la confianza en el código y proporcionaría insights adicionales sobre los mecanismos subyacentes.



% Una línea de trabajo futuro relevante para la arquitectura del framework es la reestructuración de la presentación de estímulos a los modelos. Durante la implementación del \acrshort{GSM}, los estímulos quedaron parcialmente desacoplados del modelo neuronal, lo cual representa un paso en la dirección correcta. Sin embargo, sería beneficioso formalizar esta separación de manera más sistemática, estableciendo una interfaz clara entre la generación de estímulos y los modelos que los procesan. Esta abstracción permitiría intercambiar diferentes tipos de estímulos (visuales, auditivos, multimodales) sin modificar la lógica interna de los modelos, facilitando tanto la experimentación con diferentes paradigmas como la eventual integración de múltiples modalidades sensoriales en un mismo marco computacional.
% TODO: Acá cita de nuevo nuestro paper de scikirneuromsi, esa idea no es tuya completamente.

Una línea de trabajo futuro relevante planteada por \citeauthor{paredes2025scikit} para la arquitectura del framework es la reestructuración de la presentación de estímulos a los modelos. En nuestro caso, durante la implementación del \acrshort{GSM}, los estímulos quedaron parcialmente desacoplados del modelo neuronal, lo cual representa un paso en la dirección correcta. Sin embargo, como señalan \citet{paredes2025scikit}, sería beneficioso formalizar esta separación de manera más sistemática, estableciendo una interfaz clara entre la generación de estímulos y los modelos que los procesan. Esta abstracción permitiría intercambiar diferentes tipos de estímulos (visuales, auditivos, multimodales) sin modificar la lógica interna de los modelos, facilitando tanto la experimentación con diferentes paradigmas como la eventual integración de múltiples modalidades sensoriales en un mismo marco computacional.


% Otra extensión natural es la implementación de otros modelos generativos dentro del módulo \texttt{generative}. Modelos como campos aleatorios de Markov, modelos de energía, o modelos generativos multimodales podrían incorporarse siguiendo la interfaz establecida por el \acrshort{GSM}. Esta diversidad de modelos generativos permitiría estudiar cómo diferentes supuestos sobre la estructura estadística de los estímulos afectan las dinámicas de muestreo neuronal y las propiedades emergentes de los circuitos.
% TODO: Y por qué es importante incorporarlos? Cómo contribuyen al campo de la integración multisensorial? Citas.

Otra extensión natural es la implementación de otros modelos generativos dentro del módulo \texttt{generative}. Modelos como campos aleatorios de Markov o modelos generativos multimodales podrían incorporarse siguiendo la interfaz establecida por el \acrshort{GSM}. Esto resulta particularmente interesante porque, como mostró \citet{echeveste2020cortical}, las propiedades emergentes del circuito (frecuencia de oscilaciones gamma, magnitud de transitorios, estructura de correlaciones) dependen críticamente de los momentos objetivo derivados del modelo generativo. Disponer de múltiples modelos generativos permitiría investigar sistemáticamente cómo cambian estas dinámicas al modificar el problema de inferencia que la red debe resolver, extendiendo el análisis más allá del \acrshort{GSM} visual hacia representaciones que incluyan estructura causal multimodal, como las propuestas por \citet{kording2007causal}.



% Por último, consideramos relevante explorar la aplicación del modelo a datos experimentales reales. El marco teórico de muestreo neuronal genera predicciones específicas sobre la relación entre variabilidad neural, oscilaciones gamma, y propiedades del estímulo que podrían contrastarse con registros electrofisiológicos de corteza visual. Esta validación experimental constituiría el test más riguroso del modelo y podría revelar tanto sus fortalezas como sus limitaciones para capturar la complejidad de los circuitos corticales reales.
% TODO: ¿Qué datos experimentales? Cita los estudios que podrías reproducir.


Por último, consideramos relevante explorar la aplicación del modelo a datos experimentales reales. El marco teórico de muestreo neuronal genera predicciones específicas sobre la relación entre variabilidad neural, oscilaciones gamma y propiedades del estímulo. Estas predicciones podrían contrastarse con registros electrofisiológicos de corteza visual como los analizados por \citet{banyai2019stimulus}, quienes estudiaron cómo las correlaciones de respuesta en \acrshort{V1} dependen de la complejidad del estímulo, o con los datos de \citet{haefner2016perceptual}, que vincularon la variabilidad neuronal con decisiones perceptuales en un marco de muestreo probabilístico. Esta validación experimental constituiría el test más exigente del modelo y podría revelar tanto las condiciones en que el paradigma de muestreo captura adecuadamente la dinámica cortical como aquellas en que resulta insuficiente.










